\documentclass[11pt]{article}
\usepackage{amsmath, amssymb, amsfonts, amsthm}
\usepackage{pgfplots}
\parindent 0px

\pgfplotsset{compat=1.18, width=12cm}

\begin{document}

\section{SDL Pixel Coordinates}

Let us consider a graph that represents the application window.\\
In our example, the screen has a 16 by 9 aspect ratio with pixel\\
dimensions of $W=1024px, H=576px$.\\[5pt]

\begin{tikzpicture}
\begin{axis}[y dir=reverse, axis x line=box, xtick pos=upper, xmin=0, xmax=1024, ymin=0, ymax=576]
\addplot[only marks] coordinates{(0, 0)};
\end{axis}
\end{tikzpicture}
\\

This is how SDL sees the window. Notice how the origin sits in the\\
upper-left corner of the screen. Y-values go higher as we travel \\
down the screen.\\

\rule{12cm}{0.4pt}
\\

Now, let us consider the application window using our 2D coordinate\\
system\\

\begin{tikzpicture}
\begin{axis}[axis x line=box, xmin=-1.778, xmax=1.778, ymin=-1, ymax=1]
\addplot[only marks] coordinates{(0, 0)};
\addplot[color=black] {0};
\draw[] (0,-1)--(0,1);
\end{axis}
\end{tikzpicture}
\\

The origin now sits in the center of the screen, and Y-values go higher\\
as we travel up the axis. Additionally, note the change in units.\\
One "screen unit" is the distance between the center and the top or\\
bottom of the screen. In this example, our screen has a $16:9$ aspect\\
ratio, which gives our X-values the following range:\\

$$\left\{ x \in \mathbb{R} \left| -\frac{16}{9} \leq x \leq \frac{16}{9} \right. \right\}$$
\\

We can replace $\frac{16}{9}$ with any other aspect ratio when working with\\
different windows.\\

\rule{12cm}{0.4pt}
\\

To translate a point in our 2D coordinate system to a point in\\
SDL Pixel Coordinates, consider the following analysis:\\

Let $\vec{P_{2D}}$ be any given point in our 2D coordinate system.\\
Let $\vec{P_{sdl}}$ be any given point in SDL pixel coordinates.\\
Let $\vec{P_{orig}}=\begin{bmatrix} X_{orig}\\  Y_{orig} \end{bmatrix}$ be the center of our application window in\\ 
SDL pixel coordinates, where\\

$$X_{orig} = \frac{W}{2}$$
and
$$Y_{orig} = \frac{H}{2}.$$
\\

The coordinate conversion between a point $\vec{P_{2D}}$ and the same point $\vec{P_{sdl}}$\\
can be described by the affine transformation

$$ \mathbf{C}\vec{P_{2D}} + \vec{P_{orig}} = \vec{P_{sdl}}$$
where
$$ \mathbf{C}= \begin{bmatrix} Y_{orig} &    0 \\
                                   0     & -Y_{orig} \end{bmatrix}$$
\\
is our coordinate conversion matrix. We can take the expanded expression,

$$ \begin{bmatrix} Y_{orig} & 0\\ 0 & -Y_{orig} \end{bmatrix} \begin{bmatrix}X_{2D}\\ Y_{2D}\end{bmatrix} +
\begin{bmatrix} X_{orig}\\ Y_{orig}\end{bmatrix} = \begin{bmatrix} X_{sdl}\\ Y_{sdl}\end{bmatrix}$$

and solve for each component,

\begin{align}
 Y_{orig} X_{2D} + X_{orig} &= X_{sdl}\\
-Y_{orig} Y_{2D} + Y_{orig} &= Y_{sdl}
\end{align}

\end{document}
